%======================================================================
\section{Perturbative expansion of the evolution operator ($\OM$) with quarks}
\label{sec:omega}
%======================================================================

In the previous section the evolution operator $\OM$ was defined implicitly through
Eq.~\eqref{eq:Omegadef}. Beyond leading order no closed form for $\OM$ is available and it must
be constructed perturbatively. Our strategy is the one of Paper~I. We first write down the most
general ansatz for $\OM$ as a fully normal ordered series in the soft gluon, quark and antiquark
creation and annihilation operators, $a^{\dg},a,b^{\dg},b,d^{\dg},d$, with unknown operator valued
coefficients, retaining only those operator structures that can contribute up to $\order{g^{2}}$.
The unknown coefficients are then fixed by two complementary sets of conditions:
\begin{itemize}
\item the unitarity requirement $\OM^{\dg}\OM=1$, imposed order by order in the coupling;
\item the requirement that $\OM$, acting on the free Fock states, reproduces the dressed states
computed directly in LCPT, i.e.\ $\OM\ket{n}=\ket{\Psi_{n}}$ for the vacuum, one, two and
three-gluon states, a single quark (antiquark), a quark--antiquark pair and a quark--gluon state.
\end{itemize}
Each coefficient, containing a definite number of annihilation and creation operators, is probed
only by the transition between the corresponding incoming and outgoing Fock sectors; acting on
these states therefore fixes the coefficients one by one. Since $\rho$ satisfies
\eqref{eq:rhoalg}, the coefficients are non-commuting operators on $\mathcal H_{\rm val}$ and
their ordering carries physical significance.

\subsection*{The ansatz}
\label{sec:ansatz}

We write
\begin{equation}
\OM=\OM_{\rm YM}+\OM_{q},
\label{eq:OmegaSplit}
\end{equation}
where $\OM_{\rm YM}$ is the pure gauge evolution operator of Paper~I, Eq.~(4.1) there, with
coefficients $N,A,B_{1,2,3},C,C_{2},D_{1,2,3},E_{1,2,3}$, now evaluated with the total charge
density \eqref{eq:rho}, and $\OM_{q}$ contains every normal ordered structure with at least one
soft quark or antiquark operator. The latter is organized by order and by the type of process,
\begin{equation}
\OM_{q}=\OM^{(1)}_{q}+\OM^{(2)}_{q\bar q}+\OM^{(2)}_{F}+\OM^{(2)}_{\times}.
\label{eq:Omqdecomp}
\end{equation}
At order $g$,
\begin{align}
\OM^{(1)}_{q}=&\int_{p,l,m}\Big[\mf A^{b\alpha\beta}_{3\,j\lambda_{2}\lambda_{1}}(p,m,l)\,
b^{\beta\dg}_{\lambda_{1}}(l)d^{\alpha\dg}_{\lambda_{2}}(m)a^{b}_{j}(p)
\nonumber\\&\hspace{1.6cm}
+\mf A^{\beta\alpha b}_{4\,\lambda_{1}\lambda_{2}j}(l,m,p)\,
a^{\dg b}_{j}(p)d^{\alpha}_{\lambda_{2}}(m)b^{\beta}_{\lambda_{1}}(l)\Big]
\nonumber\\
+&\int_{p,l,l'}\Big[\mf F^{\beta'\beta b}_{1\,\lambda'\lambda j}(l',l,p)\,
a^{\dg b}_{j}(p)b^{\beta\dg}_{\lambda}(l)b^{\beta'}_{\lambda'}(l')
\nonumber\\&\hspace{1.6cm}
+\mf F^{b\beta\beta'}_{2\,j\lambda\lambda'}(p,l,l')\,
b^{\beta'\dg}_{\lambda'}(l')b^{\beta}_{\lambda}(l)a^{b}_{j}(p)\Big]
\nonumber\\
+&\int_{p,m,m'}\Big[\mf G^{\alpha'\alpha b}_{1\,\lambda'\lambda j}(m',m,p)\,
a^{\dg b}_{j}(p)d^{\alpha\dg}_{\lambda}(m)d^{\alpha'}_{\lambda'}(m')
\nonumber\\&\hspace{1.6cm}
+\mf G^{b\alpha\alpha'}_{2\,j\lambda\lambda'}(p,m,m')\,
d^{\alpha'\dg}_{\lambda'}(m')d^{\alpha}_{\lambda}(m)a^{b}_{j}(p)\Big].
\label{eq:Om1q}
\end{align}
$\mf A_{3}$ splits a soft gluon into a soft pair and $\mf A_{4}$ merges a soft pair into a soft gluon;
$\mf F_{1}$, $\mf G_{1}$ ($\mf F_{2}$, $\mf G_{2}$) describe the emission (absorption) of a soft gluon
by a soft quark and by a soft antiquark. At order $g^{2}$, the structures that create or
annihilate soft pairs are
{\small
\begin{align}
\OM^{(2)}_{q\bar q}=&\int_{l,m}\Big[\mf A^{\alpha\beta}_{1\,\lambda_{2}\lambda_{1}}(m,l)\,
b^{\beta\dg}_{\lambda_{1}}(l)d^{\alpha\dg}_{\lambda_{2}}(m)
+\mf A^{\beta\alpha}_{2\,\lambda_{1}\lambda_{2}}(l,m)\,d^{\alpha}_{\lambda_{2}}(m)b^{\beta}_{\lambda_{1}}(l)\Big]
\nonumber\\
+&\int_{p,q,l,m}\Big[\mf B^{cb\alpha\beta}_{1\,kj\lambda_{2}\lambda_{1}}(q,p,m,l)\,
b^{\beta\dg}_{\lambda_{1}}(l)d^{\alpha\dg}_{\lambda_{2}}(m)a^{b}_{j}(p)a^{c}_{k}(q)
\nonumber\\
&\qquad+\mf B^{\beta\alpha cb}_{2\,\lambda_{1}\lambda_{2}kj}(l,m,q,p)\,
a^{\dg b}_{j}(p)a^{\dg c}_{k}(q)d^{\alpha}_{\lambda_{2}}(m)b^{\beta}_{\lambda_{1}}(l)
\nonumber\\
&\qquad+\mf B^{\alpha\beta bc}_{3\,\lambda_{2}\lambda_{1}jk}(m,l,p,q)\,
a^{\dg c}_{k}(q)a^{b}_{j}(p)b^{\beta\dg}_{\lambda_{1}}(l)d^{\alpha\dg}_{\lambda_{2}}(m)
\nonumber\\
&\qquad+\mf B^{cb\beta\alpha}_{4\,kj\lambda_{1}\lambda_{2}}(q,p,l,m)\,
d^{\alpha}_{\lambda_{2}}(m)b^{\beta}_{\lambda_{1}}(l)a^{\dg b}_{j}(p)a^{c}_{k}(q)\Big]
\nonumber\\
+&\int_{p,q,r,s,l,m}\Big[
\mf C^{\beta\alpha bcde}_{1\,\lambda_{1}\lambda_{2}jkon}(l,m,p,q,r,s)\,
a^{\dg e}_{n}(s)a^{\dg d}_{o}(r)a^{\dg c}_{k}(q)a^{b}_{j}(p)d^{\alpha}_{\lambda_{2}}(m)b^{\beta}_{\lambda_{1}}(l)
\nonumber\\
&\qquad+\mf C^{\beta\alpha bcde}_{2\,\lambda_{1}\lambda_{2}jkon}(l,m,p,q,r,s)\,
a^{\dg e}_{n}(s)a^{\dg d}_{o}(r)a^{c}_{k}(q)a^{b}_{j}(p)d^{\alpha}_{\lambda_{2}}(m)b^{\beta}_{\lambda_{1}}(l)
\nonumber\\
&\qquad+\mf C^{bcde\alpha\beta}_{3\,jkon\lambda_{2}\lambda_{1}}(p,q,r,s,m,l)\,
b^{\beta\dg}_{\lambda_{1}}(l)d^{\alpha\dg}_{\lambda_{2}}(m)a^{\dg e}_{n}(s)a^{d}_{o}(r)a^{c}_{k}(q)a^{b}_{j}(p)
\nonumber\\
&\qquad+\mf C^{bcde\alpha\beta}_{4\,jkon\lambda_{2}\lambda_{1}}(p,q,r,s,m,l)\,
b^{\beta\dg}_{\lambda_{1}}(l)d^{\alpha\dg}_{\lambda_{2}}(m)a^{\dg e}_{n}(s)a^{\dg d}_{o}(r)a^{c}_{k}(q)a^{b}_{j}(p)\Big]
\nonumber\\
+&\int_{p,q,l,m,l',m'}\Big[
\mf D^{cb\alpha'\beta'\alpha\beta}_{1\,kj\lambda_{2}'\lambda_{1}'\lambda_{2}\lambda_{1}}(q,p,m',l',m,l)
\nonumber\\
&\qquad\quad\times b^{\beta\dg}_{\lambda_{1}}(l)d^{\alpha\dg}_{\lambda_{2}}(m)b^{\beta'\dg}_{\lambda_{1}'}(l')d^{\alpha'\dg}_{\lambda_{2}'}(m')a^{b}_{j}(p)a^{c}_{k}(q)
\nonumber\\
&\qquad+\mf D^{\beta\alpha\beta'\alpha'cb}_{2\,\lambda_{1}\lambda_{2}\lambda_{1}'\lambda_{2}'kj}(l,m,l',m',q,p)
\nonumber\\
&\qquad\quad\times a^{\dg b}_{j}(p)a^{\dg c}_{k}(q)d^{\alpha'}_{\lambda_{2}'}(m')b^{\beta'}_{\lambda_{1}'}(l')d^{\alpha}_{\lambda_{2}}(m)b^{\beta}_{\lambda_{1}}(l)\Big]
\nonumber\\
+&\int_{l,m,l',m'}
\mf D^{\beta'\alpha'\alpha\beta}_{3\,\lambda_{1}'\lambda_{2}'\lambda_{2}\lambda_{1}}(l',m',m,l)\,
b^{\beta\dg}_{\lambda_{1}}(l)d^{\alpha\dg}_{\lambda_{2}}(m)d^{\alpha'}_{\lambda_{2}'}(m')b^{\beta'}_{\lambda_{1}'}(l') .
\label{eq:Om2qq}
\end{align}}
The labels of each coefficient are listed in the order in which the operators appear when the
string is read from right to left. Every pair is annihilated in the order $d\,b$ and created in
the order $b^{\dg}d^{\dg}$; with this convention all the matrix elements of $\OM$ between states
containing one soft pair, defined in appendix~\ref{app:states}, carry a positive sign. The
structures that describe soft quark (antiquark) lines are
{\small
\begin{align}
\OM^{(2)}_{F}=&\int_{l,l'}\mf P^{\beta'\beta}_{\lambda'\lambda}(l',l)\,
b^{\beta\dg}_{\lambda}(l)b^{\beta'}_{\lambda'}(l')
\nonumber\\
&+\int_{q,r,l,l'}\Big[\mf Q^{\beta'\beta cd}_{1\,\lambda'\lambda kn}(l',l,q,r)\,
a^{\dg d}_{n}(r)a^{\dg c}_{k}(q)b^{\beta\dg}_{\lambda}(l)b^{\beta'}_{\lambda'}(l')
\nonumber\\
&\qquad+\mf Q^{\beta'c\beta d}_{2\,\lambda'k\lambda n}(l',q,l,r)\,
a^{\dg d}_{n}(r)b^{\beta\dg}_{\lambda}(l)a^{c}_{k}(q)b^{\beta'}_{\lambda'}(l')
\nonumber\\
&\qquad+\mf Q^{dc\beta\beta'}_{3\,nk\lambda\lambda'}(r,q,l,l')\,
b^{\beta'\dg}_{\lambda'}(l')b^{\beta}_{\lambda}(l)a^{c}_{k}(q)a^{d}_{n}(r)\Big]
\nonumber\\
&+\int_{l,l',l_{2},m}\Big[\mf K^{\beta'\beta\beta_{2}\alpha}_{1\,\lambda'\lambda\lambda_{2}\lambda_{m}}(l',l,l_{2},m)\,
b^{\beta\dg}_{\lambda}(l)b^{\beta_{2}\dg}_{\lambda_{2}}(l_{2})d^{\alpha\dg}_{\lambda_{m}}(m)b^{\beta'}_{\lambda'}(l')
\nonumber\\
&\qquad+\mf K^{\alpha\beta_{2}\beta\beta'}_{2\,\lambda_{m}\lambda_{2}\lambda\lambda'}(m,l_{2},l,l')\,
b^{\beta'\dg}_{\lambda'}(l')b^{\beta}_{\lambda}(l)b^{\beta_{2}}_{\lambda_{2}}(l_{2})d^{\alpha}_{\lambda_{m}}(m)\Big]
+\big(q\to\bar q\big),
\label{eq:OmF}
\end{align}}
where $(q\to\bar q)$ denotes the analogous antiquark structures, with coefficients
$\bar{\mf P},\bar{\mf Q}_{1,2,3},\bar{\mf K}_{1,2}$, obtained by $b\leftrightarrow d$. $\mf P$
($\bar{\mf P}$) describes the eikonal rescattering of a soft quark (antiquark) on the valence
charge; $\mf Q_{1}$ the emission of two soft gluons by a soft quark, $\mf Q_{2}$ quark--gluon
scattering and $\mf Q_{3}$ the absorption of two gluons; $\mf K_{1}$ the production of a soft pair
off a soft quark and $\mf K_{2}$ the reverse process. Finally, $\OM^{(2)}_{\times}$ collects the
remaining products of two first order structures that act on disjoint sets of particles,
\begin{equation}
\OM^{(2)}_{\times}=\frac12:\!\OM^{(1)}\OM^{(1)}\!:\Big|_{\text{no contraction, not in }\eqref{eq:Om2qq},\eqref{eq:OmF}},
\qquad \OM^{(1)}\equiv\OM^{(1)}_{\rm YM}+\OM^{(1)}_{q},
\label{eq:Omx}
\end{equation}
where $\OM^{(1)}_{\rm YM}$ contains the $A$ and $C$ terms of Paper~I and $:\;:$ denotes normal
ordering. These six-operator structures (e.g.\ $\mf A_{3}\,\mf F_{1}$ or $\mf F_{1}\,C$) are fixed
completely by \eqref{eq:factor}, as shown in section~\ref{sec:coeff}; we list them separately
only to keep \eqref{eq:Om2qq} and \eqref{eq:OmF} readable. The structures $\mf C_{1,\dots,4}$ and
$\mf D_{1,2}$ in \eqref{eq:Om2qq} are of the same type, but we keep them explicitly since they act
on the states considered below. Normal ordered structures with more operators than those displayed
contribute beyond $\order{g^{2}}$ and are omitted. The coefficients $\mf A_{3}$, $\mf A_{4}$,
$\mf F_{1,2}$, $\mf G_{1,2}$ are of order $g$; all others are of order $g^{2}$. The identity
component of $\OM$ is the normalization $N$ of Paper~I; since soft pairs are not created out of
the vacuum at order $g$, $N$ receives no quark contribution at $\order{g^{2}}$.

\subsection*{Unitarity constraints}
\label{sec:unitarity}

Expanding $\OM=1+g\OM^{(1)}+g^{2}\OM^{(2)}+\order{g^{3}}$, the unitarity condition gives, as in
Paper~I,
\begin{equation}
\OM^{\dg(1)}=-\OM^{(1)},\qquad
\OM^{(2)}+\OM^{\dg(2)}=-\OM^{\dg(1)}\OM^{(1)}=\OM^{(1)}\OM^{(1)} .
\label{eq:U12}
\end{equation}
In the quark sector the first condition gives
\begin{align}
\mf A^{\beta\alpha b}_{4\,\lambda_{1}\lambda_{2}j}(l,m,p)&=-\big[\mf A^{b\alpha\beta}_{3\,j\lambda_{2}\lambda_{1}}(p,m,l)\big]^{\dg},
\nonumber\\
\mf F^{b\beta\beta'}_{2\,j\lambda\lambda'}(p,l,l')&=-\big[\mf F^{\beta'\beta b}_{1\,\lambda'\lambda j}(l',l,p)\big]^{\dg},
\qquad
\mf G^{b\alpha\alpha'}_{2\,j\lambda\lambda'}(p,m,m')=-\big[\mf G^{\alpha'\alpha b}_{1\,\lambda'\lambda j}(m',m,p)\big]^{\dg}.
\label{eq:U1}
\end{align}
From here onward we write $\mf F\equiv\mf F_{1}$, $\mf G\equiv\mf G_{1}$. The second condition,
projected onto the independent normal ordered structures, gives the following relations. To
lighten the notation we suppress the colour, polarization and helicity labels: each coefficient
carries the labels of its momentum arguments, in the order of Eqs.~\eqref{eq:Om1q}--\eqref{eq:OmF},
and repeated internal momenta are integrated with $\int$.
\begin{align}
\mf A_{2}(l,m)&=-\mf A_{1}^{\dg}(m,l)+\int_{p}A^{\dg}(p)\,\mf A_{3}^{\dg}(p,m,l),
\label{eq:UA2}\\[3pt]
\mf B_{2}(l,m,q,p)&=-\mf B_{1}^{\dg}(q,p,m,l)-\mf A_{3}^{\dg}(p,m,l)A(q)-A(p)\mf A_{3}^{\dg}(q,m,l)
\nonumber\\
&\quad-\int_{u}C(u,q,p)\,\mf A_{3}^{\dg}(u,m,l)
\nonumber\\
&\quad-\int_{l_{1}}\mf A_{3}^{\dg}(q,m,l_{1})\,\mf F(l,l_{1},p)
-\int_{m_{1}}\mf A_{3}^{\dg}(q,m_{1},l)\,\mf G(m,m_{1},p),
\label{eq:UB2}\\[3pt]
\mf B_{4}(q,p,l,m)&=-\mf B_{3}^{\dg}(m,l,p,q)+\mf A_{3}^{\dg}(p,m,l)A^{\dg}(q)+A^{\dg}(q)\mf A_{3}^{\dg}(p,m,l)
\nonumber\\
&\quad+2\int_{u}C^{\dg}(p,q,u)\,\mf A_{3}^{\dg}(u,m,l)
\nonumber\\
&\quad+\int_{l_{1}}\mf A_{3}^{\dg}(p,m,l_{1})\,\mf F^{\dg}(l_{1},l,q)
+\int_{m_{1}}\mf A_{3}^{\dg}(p,m_{1},l)\,\mf G^{\dg}(m_{1},m,q),
\label{eq:UB4}\\[3pt]
\mf C_{1}(l,m,p,q,r,s)&=-\mf C_{3}^{\dg}(p,q,r,s,m,l)
-\mathcal S_{3}\Big\{\mf A_{3}^{\dg}(s,m,l),\,C(p,q,r)\Big\},
\label{eq:UC1}\\[3pt]
\mf C_{2}(l,m,p,q,r,s)&=-\mf C_{4}^{\dg}(p,q,r,s,m,l)
+\mathcal S_{2}\Big\{\mf A_{3}^{\dg}(s,m,l),\,C^{\dg}(r,q,p)\Big\},
\label{eq:UC2}\\[3pt]
\mf D_{2}(l,m,l',m',q,p)&=-\mf D_{1}^{\dg}(q,p,m',l',m,l)
\nonumber\\&\quad
+\mf A_{3}^{\dg}(p,m,l)\mf A_{3}^{\dg}(q,m',l')+\mf A_{3}^{\dg}(q,m,l)\mf A_{3}^{\dg}(p,m',l'),
\label{eq:UD2}\\[3pt]
\mf Q_{3}(r,q,l,l')&=-\mf Q_{1}^{\dg}(l',l,q,r)+A^{\dg}(q)\mf F^{\dg}(l',l,r)
\nonumber\\
&\quad+\mf F^{\dg}(l',l,r)A^{\dg}(q)
\nonumber\\
&\quad+\int_{l_{1}}\mf F^{\dg}(l_{1},l,r)\,\mf F^{\dg}(l',l_{1},q)
+\int_{s}\mf F^{\dg}(l',l,s)\,C^{\dg}(s,q,r)+(q\leftrightarrow r),
\label{eq:UQ3}\\[3pt]
\mf K_{2}(m,l_{2},l,l')&=-\mf K_{1}^{\dg}(l',l,l_{2},m)
+\Big[\int_{k}\mf F^{\dg}(l',l,k)\,\mf A_{3}^{\dg}(k,m,l_{2})-\big(l\leftrightarrow l_{2}\big)\Big].
\label{eq:UK2}
\end{align}
Here $\mathcal S_{3}$ symmetrizes over the three gluon labels created by $\mf C_{3}^{\dg}$ and
$\mathcal S_{2}$ over the two created by $\mf C_{4}^{\dg}$, $\{X,Y\}\equiv XY+YX$, $A$ and $C$ are the
Weizs\"acker--Williams and three-gluon coefficients of Paper~I, and the antisymmetrization in
\eqref{eq:UK2} acts on all the labels of the two identical quarks. The antiquark relations follow
by $\mf F\to\mf G$, $b\to d$. The structures $\mf B_{2}$, $\mf B_{4}$ and $\mf Q_{3}$ receive
contributions from the emission and absorption coefficients $\mf F$ and $\mf G$: they correspond to
the $t$- and $u$-channel quark exchange in $gg\to q\bar q$ and to the sequential emission of two
gluons by a quark, and they are absent if the ansatz is restricted to the pair structures.

The coefficients that preserve the particle content, $\mf D_{3}$, $\mf P$, $\bar{\mf P}$ and
$\mf Q_{2}$, are fixed completely by matching to the LCWFs below. Their unitarity relations are
\begin{align}
\mf D_{3}(l',m',m,l)+\mf D_{3}^{\dg}(l,m,m',l')&=
-\int_{k}\mf A_{3}(k,m,l)\,\mf A_{3}^{\dg}(k,m',l')
\nonumber\\
&\quad-\int_{p}\Big[\mf G^{\dg}(m,m',p)\,\mf F(l',l,p)+\mf F^{\dg}(l,l',p)\,\mf G(m',m,p)\Big],
\label{eq:UD3}\\[3pt]
\mf P(l',l)+\mf P^{\dg}(l,l')&=-\int_{p}\Big[A^{\dg}(p)\,\mf F(l',l,p)+\mf F^{\dg}(l,l',p)\,A(p)\Big]
\nonumber\\
&\quad-\int_{p,l_{1}}\mf F^{\dg}(l,l_{1},p)\,\mf F(l',l_{1},p),
\label{eq:UP}\\[3pt]
\mf Q_{2}(f\leftarrow i)+\mf Q_{2}^{\dg}(i\leftarrow f)&=-(2\pi)^{6}\sum_{m}\big[\mu_{1}(m\leftarrow f)\big]^{\dg}\mu_{1}(m\leftarrow i),
\label{eq:UQ2}
\end{align}
and they are therefore nontrivial checks of the calculation. In \eqref{eq:UQ2} the intermediate
states are $q$, $qgg$, $qq\bar q$ and, for the $\rho$ dependent part, $qgg$ with one gluon
exchanged with the valence; each term is the product of two of the first order coefficients
$A,C,\mf A_{3},\mf F$. All of them are verified intermediate state by intermediate state with the
identity \eqref{eq:unitid}. The relation for $\mf D_{3}$ requires the $t$-channel exchange of a
gluon between the quark and the antiquark, and the relation for $\mf P$ requires the exchange of a
gluon between the soft quark and the valence charge; both involve precisely the coefficients
$\mf F$ and $\mf G$. Finally, the pure gauge relation for the coefficient $B_{3}$ of Paper~I (the $a^{\dg}a$
structure) acquires the quark loop term $-\int_{l,m}\mf A_{3}^{\dg}(k',m,l)\mf A_{3}(k,m,l)$ on its
right-hand side; it affects only the momentum diagonal part of $B_{3}$, discussed in
section~\ref{sec:wf-g}.

\subsection*{Action of $\OM$ on Fock states}
\label{sec:action}

Since $\OM$ is normal ordered, each structure contributes to the action on a given Fock state in
a simple way: every annihilation operator either contracts with an incoming particle or
annihilates the vacuum, while the creation operators produce the particles of the outgoing
component. If a structure contains $n$ annihilation and $M$ creation operators and acts on an
$n$-particle state, it produces
\begin{equation}
\OM\ket{n}\supset\frac{1}{(2\pi)^{3(n+M)/2}}\int\!d^{3}(M\ \text{momenta})\;K\;\ket{M},
\qquad\text{i.e.}\qquad K=(2\pi)^{3(n+M)/2}\,\psi,
\label{eq:rule}
\end{equation}
where $K$ is the coefficient and $\psi$ the coefficient of the outgoing state in the LCWF
\eqref{eq:fock}. We list only the terms that involve the quark sector and carry new information
through $\order{g^{2}}$; the purely gluonic terms are those of Paper~I.

\paragraph{Vacuum.}
\begin{equation}
\OM\vac\supset\frac{1}{(2\pi)^{3}}\int\!d^{3}l\,d^{3}m\;\mf A^{\alpha\beta}_{1\,\lambda_{2}\lambda_{1}}(m,l)\,
\ket{\qk{\beta}{\lambda_{1}}{l}\,\aqk{\alpha}{\lambda_{2}}{m}} .
\label{eq:Omvac}
\end{equation}

\paragraph{One gluon.}
\begin{align}
\OM\ket{\gl{a}{i}{k}}\supset\;&\frac{1}{(2\pi)^{9/2}}\int\!d^{3}l\,d^{3}m\;
\mf A^{a\alpha\beta}_{3\,i\lambda_{2}\lambda_{1}}(k,m,l)\,\ket{\qk{\beta}{\lambda_{1}}{l}\aqk{\alpha}{\lambda_{2}}{m}}
\nonumber\\
&+\frac{1}{(2\pi)^{6}}\int\!d^{3}l\,d^{3}m\,d^{3}q\;\mf B^{\alpha\beta ac}_{3\,\lambda_{2}\lambda_{1}ik}(m,l,k,q)\,
\ket{\qk{\beta}{\lambda_{1}}{l}\aqk{\alpha}{\lambda_{2}}{m}\gl{c}{k}{q}}
\nonumber\\
&+\frac{1}{(2\pi)^{3}}\int\!d^{3}l\,d^{3}m\;\mf A_{1}(m,l)\,\ket{\gl{a}{i}{k}\,q\bar q}
\nonumber\\
&+\frac{1}{(2\pi)^{3}}\int\!d^{3}p\;B_{3}(k,p)\Big|_{\rm quark\ loop}\ket{\gl{a}{i}{p}} .
\label{eq:Omg}
\end{align}

\paragraph{Two gluons.}
\begin{align}
\OM\ket{\gl{a}{i}{k}\gl{b}{j}{p}}\supset\;&\frac{1}{(2\pi)^{6}}\int\!d^{3}l\,d^{3}m\;
\ket{\qk{\beta}{\lambda_{1}}{l}\aqk{\alpha}{\lambda_{2}}{m}}
\nonumber\\&\hspace{1.2cm}\times\Big[\mf B^{ba\alpha\beta}_{1\,ji\lambda_{2}\lambda_{1}}(p,k,m,l)+\mf B^{ab\alpha\beta}_{1\,ij\lambda_{2}\lambda_{1}}(k,p,m,l)\Big]
\nonumber\\
&+\frac{1}{(2\pi)^{9}}\int\!d^{3}r\,d^{3}s\,d^{3}l\,d^{3}m\;
\Big[\mf C_{4}(k,p,r,s,m,l)+\mf C_{4}(p,k,r,s,m,l)\Big]
\nonumber\\&\hspace{1.2cm}\times\ket{q\bar q\,\gl{d}{o}{r}\gl{e}{n}{s}}
\nonumber\\
&+\frac{1}{(2\pi)^{9}}\int\!d^{3}l\,d^{3}m\,d^{3}l'\,d^{3}m'\;
\Big[\mf D_{1}(p,k,m',l',m,l)+\mf D_{1}(k,p,m',l',m,l)\Big]
\nonumber\\&\hspace{1.2cm}\times\ket{q\bar q\,q\bar q}
\nonumber\\
&+\frac{1}{(2\pi)^{9/2}}\int\!d^{3}l\,d^{3}m\;\Big[\mf A_{3}(k,m,l)\ket{\gl{b}{j}{p}\,q\bar q}
\nonumber\\&\hspace{3.6cm}
+\mf A_{3}(p,m,l)\ket{\gl{a}{i}{k}\,q\bar q}\Big]+\dots
\label{eq:Omgg}
\end{align}
where the dots denote the disconnected terms $\mf A_{1}\ket{gg\,q\bar q}$ and
$\mf B_{3}\ket{g\,q\bar q\,g}$ with a spectator gluon.

\paragraph{Three gluons.}
\begin{equation}
\OM\ket{\gl{a}{i}{k}\gl{b}{j}{p}\gl{c}{k}{q}}\supset\frac{1}{(2\pi)^{9}}\int\!d^{3}s\,d^{3}l\,d^{3}m\;
\sum_{3\ \rm assign.}\mf C_{3}(k,p,q,s,m,l)\,\ket{q\bar q\,\gl{e}{n}{s}}+\dots
\label{eq:Omggg}
\end{equation}

\paragraph{One quark.}
\begin{align}
\OM\ket{\qk{\beta'}{\lambda'}{l'}}\supset\;&\frac{1}{(2\pi)^{3}}\int\!d^{3}l\;
\mf P^{\beta'\beta}_{\lambda'\lambda}(l',l)\ket{\qk{\beta}{\lambda}{l}}
\nonumber\\
&+\frac{1}{(2\pi)^{9/2}}\int\!d^{3}l\,d^{3}p\;\mf F^{\beta'\beta b}_{\lambda'\lambda j}(l',l,p)
\ket{\qk{\beta}{\lambda}{l}\gl{b}{j}{p}}
\nonumber\\
&+\frac{1}{(2\pi)^{6}}\int\!d^{3}l\,d^{3}q\,d^{3}r\;\mf Q_{1}(l',l,q,r)\ket{\qk{\beta}{\lambda}{l}\gl{c}{k}{q}\gl{d}{n}{r}}
\nonumber\\
&+\frac{1}{(2\pi)^{6}}\int\!d^{3}l\,d^{3}l_{2}\,d^{3}m\;\mf K_{1}(l',l,l_{2},m)\ket{q\,q\bar q}
\nonumber\\
&+\frac{1}{(2\pi)^{3/2}}\int\!d^{3}p\;A(p)\ket{\qk{\beta'}{\lambda'}{l'}\gl{b}{j}{p}}+\dots
\label{eq:Omq}
\end{align}
and analogously for an antiquark with $\bar{\mf P},\mf G,\bar{\mf Q}_{1},\bar{\mf K}_{1}$.

\paragraph{Quark--antiquark pair.}
\begin{align}
\OM\ket{\qk{\beta'}{\lambda_{1}'}{l'}\aqk{\alpha'}{\lambda_{2}'}{m'}}\supset\;&
\frac{1}{(2\pi)^{3}}\,\mf A_{2}(l',m')\vac
\nonumber\\
&+\frac{1}{(2\pi)^{9/2}}\int\!d^{3}k\;\mf A_{4}(l',m',k)\ket{\gl{a}{i}{k}}
\nonumber\\
&+\frac{1}{(2\pi)^{6}}\int\!d^{3}l\,d^{3}m\;\mf D_{3}(l',m',m,l)\ket{\qk{\beta}{\lambda_{1}}{l}\aqk{\alpha}{\lambda_{2}}{m}}
\nonumber\\
&+\frac{1}{(2\pi)^{3}}\int\!d^{3}l\;\mf P(l',l)\ket{\qk{\beta}{\lambda_{1}}{l}\aqk{\alpha'}{\lambda_{2}'}{m'}}
\nonumber\\
&+\frac{1}{(2\pi)^{3}}\int\!d^{3}m\;\bar{\mf P}(m',m)\ket{\qk{\beta'}{\lambda_{1}'}{l'}\aqk{\alpha}{\lambda_{2}}{m}}
\nonumber\\
&+\frac{1}{(2\pi)^{6}}\int\!d^{3}p\,d^{3}q\;\mf B_{2}(l',m',q,p)\ket{\gl{b}{j}{p}\gl{c}{k}{q}}
\nonumber\\
&+\frac{1}{(2\pi)^{9/2}}\int\!d^{3}l\,d^{3}p\;\mf F(l',l,p)\ket{q\,\bar q\,g}
\nonumber\\
&+\frac{1}{(2\pi)^{9/2}}\int\!d^{3}m\,d^{3}p\;\mf G(m',m,p)\ket{q\,\bar q\,g}+\dots
\label{eq:Omqq}
\end{align}

\paragraph{Quark--gluon state.}
\begin{align}
\OM\ket{\qk{\beta'}{\lambda'}{l'}\gl{c}{k}{q}}\supset\;&\frac{1}{(2\pi)^{6}}\int\!d^{3}l\,d^{3}r\;
\mf Q_{2}(l',q,l,r)\ket{\qk{\beta}{\lambda}{l}\gl{d}{n}{r}}
\nonumber\\
&+\frac{1}{(2\pi)^{9/2}}\int\!d^{3}l\;\mf F_{2}(q,l',l)\ket{\qk{\beta}{\lambda}{l}}+\dots
\label{eq:Omqg}
\end{align}

The coefficients of real transitions, $\mf A_{1}$, $\mf A_{3}$, $\mf F$, $\mf G$, $\mf B_{1}$,
$\mf B_{3}$, $\mf C_{3}$, $\mf C_{4}$, $\mf D_{1}$, $\mf Q_{1}$, $\mf K_{1}$, and the particle number
preserving coefficients $\mf D_{3}$, $\mf P$, $\bar{\mf P}$, $\mf Q_{2}$ are determined by matching
these expressions to the LCWFs computed in the following subsections. The remaining coefficients,
$\mf A_{2}$, $\mf A_{4}$, $\mf F_{2}$, $\mf G_{2}$, $\mf B_{2}$, $\mf B_{4}$, $\mf C_{1}$, $\mf C_{2}$,
$\mf D_{2}$, $\mf Q_{3}$, $\mf K_{2}$, follow from the unitarity relations.\footnote{For
$\mf A_{4}$, $\mf A_{2}$ and $\mf B_{2}$ we have also verified the unitarity relations by a direct
LCPT computation of the transitions $q\bar q\to g$, $q\bar q\to0$ and $q\bar q\to gg$; see
section~\ref{sec:wf-qq}.}
